% ============================================================
% Project Proposal -- 605.745 Reasoning Under Uncertainty
% ============================================================
% Build with:  .\build.ps1
% ============================================================
\documentclass[11pt]{article}
\usepackage{jhu}

% ------------------------------------------------------------
% Fill these in
% ------------------------------------------------------------
\renewcommand{\jhuTitle}{Particle Filter vs. Extended Kalman Filter Localization for a Mobile Manipulator}
\renewcommand{\jhuDocType}{Project Proposal}
\renewcommand{\jhuAuthor}{Jacob Taylor Cassady}
\renewcommand{\jhuCourse}{605.745: Reasoning Under Uncertainty}
\renewcommand{\jhuRunningHead}{605.745 Final Project Proposal}

\begin{document}

\jhutitlepage
% \jhucontents   % omitted: the proposal is one page, a contents page adds nothing

\section{Problem Definition}
The broad goal of this project is to develop and test localization algorithms for a mobile robot fitted with front- and rear-facing 2D LIDARs.
Localization is the process of determining where an agent is in a given map.
This capability is a prerequisite for a robotic system's planning pipeline.
This project will perform a study comparing localization methods in a Gazebo/ROS~2 simulation environment leveraging \texttt{devol}, an existing repository on the author's public GitHub, which includes models for a Clearpath Husky A200 mobile platform carrying a UR3e arm and a Robotiq 2F-85 gripper.
The robot will be in an open-source factory environment with a known occupancy map of obstacles.
Its only information about its own pose comes from two imperfect sources: wheel odometry and LIDAR scans.

The wheel odometry comes from Gazebo's built-in odometry plugin, which models drift and slip, so pose uncertainty accumulates without bound.
The LIDAR range scans come from the ROS~2 LIDAR plugin; they are noisy and ambiguous in structured indoor environments, where one corridor segment of the factory produces nearly the same scan signature as other, distinct poses.
Localization in this environment is therefore a genuinely uncertain inference problem. 
The simulator's true model pose will be bridged into ROS~2 and used for two purposes: it drives the waypoint controller, so estimator error never alters the trajectory being scored, and it serves as ground truth for evaluation.

\section{Framework Selection and Justification}

The candidate framework is a \textbf{particle filter} (sequential Monte Carlo, bootstrap/SIR form)~\cite{gordon1993novel, arulampalam2002tutorial} implemented from scratch, following the Monte Carlo Localization formulation~\cite{dellaert1999monte, fox1999monte, thrun2005probabilistic}: the posterior over pose $(x, y, \theta)$ is a weighted sample set, propagated by an odometry motion model and re-weighted by a likelihood-field LIDAR measurement model against the known map, with low-variance resampling.
The named alternative is an \textbf{extended Kalman filter} localizer~\cite{thrun2005probabilistic, moore2014generalized}, adapted from Introduction to Robotics classwork, that fuses the same odometry with a scan-matched pose observation.
The two make opposite modeling commitments.
The EKF assumes the posterior is a single Gaussian and linearizes the motion and measurement models about the current mean.
The particle filter makes no distributional assumption and can represent a multimodal posterior which arises when a scan is consistent with several poses, when the initial pose is unknown, or when the robot is displaced without its knowledge.
The concrete reason for preferring the particle filter in this domain is therefore a testable claim: it should hold its accuracy where the posterior is multimodal, while the EKF commits to one mode and cannot recover.
Conversely, in benign runs with a known starting pose the EKF is expected to match or outperform the particle filter at a fraction of the computational cost, since its Gaussian assumption holds and it carries no sampling noise.

\section{Evaluation Plan}

Three estimators run on identical inputs: (1) dead reckoning from odometry alone, (2) an EKF, and (3) a particle filter. The EKF and particle filter receive the same odometry and LIDAR streams. Each estimator is evaluated over the same three-waypoint factory route for 20 trials with distinct random seeds against simulator ground truth using position RMSE, heading RMSE, waypoint position error, and per-step compute time.

Two adversarial tests measure recovery time to below 0.25~m position error following (a) global localization from a uniform prior and (b) mid-run kidnapping via Gazebo teleportation. 
Sensitivity analyses vary particle count, $N \in \{100, 500, 2000, 5000\}$, and injected odometry and LIDAR noise to evaluate the accuracy-compute trade-off and robustness to degraded evidence. 
Results are reported as mean $\pm$ 95\% confidence interval across trials.

The primary hypothesis is that the particle filter recovers in both adversarial scenarios where the EKF does not, and that the EKF matches or exceeds the particle filter's nominal-route accuracy at lower per-step compute time.

\newpage
\section{Acknowledgements}

An initial draft of this proposal was generated with Anthropic's Claude and substantially rewritten by the author; Claude was also used to review the revised draft against the class rubric.

\bibliography{ref}

\end{document}
